\chapter{FlashAttention 算法}
\label{chap:flash-attention}

\section{标准 Attention 的 IO 瓶颈}
\label{sec:standard-attention-io-bottleneck}

上一章讨论了 PagedAttention。PagedAttention 解决的是在线 serving 中 KV Cache 的\textbf{内存管理问题}：请求长度不一、动态加入退出、KV Cache 增长和释放，都会造成碎片与 over-reservation。PagedAttention 通过 block table 让逻辑连续的 KV Cache 可以物理离散存储。

本章讨论另一个方向：\textbf{Attention 计算本身的 IO 瓶颈}。这就是 FlashAttention 要解决的问题。

标准 Attention 的数学公式非常简单：

\[
\operatorname{Attention}(\mat{Q},\mat{K},\mat{V})
\softmax
\left(
\frac{\mat{Q}\mat{K}^{\mathsf{T}}}{\sqrt{d_h}}
\right)
\mat{V}.
\]

如果只从数学上看，计算过程似乎只是三步：

\[
\mat{S}=\frac{\mat{Q}\mat{K}^{\mathsf{T}}}{\sqrt{d_h}},
\]
\[
\mat{P}=\softmax(\mat{S}),
\]
\[
\mat{O}=\mat{P}\mat{V}.
\]

但在 GPU 上，性能不只取决于 FLOPs，还取决于数据如何在 HBM、L2、shared memory、register 之间移动。标准 Attention 最大的问题，是它通常会显式物化巨大的中间矩阵：

\[
\mat{S}\in\mathbb{R}^{S\times S},
\quad
\mat{P}\in\mathbb{R}^{S\times S}.
\]

其中 $S$ 是序列长度。随着上下文长度增长，$S^2$ 级别的中间矩阵会带来巨大的 HBM 读写压力。FlashAttention 的核心不是改变 Attention 的数学结果，而是改变计算路径：\textbf{不把完整 attention matrix 写入 HBM，而是在片上 SRAM/register 中分块完成 softmax 和 weighted sum}。

\subsection{$QK^{\mathsf{T}}$ 的中间矩阵}
\label{subsec:qk-intermediate-matrix}

考虑单个 batch、单个 attention head。设：

\[
\mat{Q},\mat{K},\mat{V}\in\mathbb{R}^{S\times d_h}.
\]

标准 Attention 第一步计算 score matrix：

\[
\mat{S}=\mat{Q}\mat{K}^{\mathsf{T}},
\]

其中：

\[
\mat{S}\in\mathbb{R}^{S\times S}.
\]

第 $i$ 行第 $j$ 列为：

\[
S_{ij}
\langle \vect{q}_i,\vect{k}_j\rangle.
\]

对于 causal attention，还需要 mask：

\[
S_{ij}=-\infty,\quad j>i.
\]

然后对每一行做 softmax：

\[
P_{ij}
\frac{\exp(S_{ij}-m_i)}
{\sum_{j'}\exp(S_{ij'}-m_i)},
\]
其中 $m_i=\max_j S_{ij}$ 用于数值稳定。

最后输出：

\[
\vect{o}_i
\sum_j P_{ij}\vect{v}_j.
\]

问题出现在 $\mat{S}$ 和 $\mat{P}$。如果 $S=4096$，则单 head 的 attention matrix 元素个数为：

\[
4096^2=16{,}777{,}216.
\]

若使用 FP16/BF16，每个元素 2 bytes，则一个矩阵约为：

\[
32\ \mathrm{MB}.
\]

如果同时保存 score、softmax probability、dropout mask 或 backward 所需中间结果，显存和 HBM traffic 会更高。对于多 head、多 batch、多层，这个开销会迅速放大。

\[
M_{\mathrm{attn\ matrix}}
B\cdot n_h\cdot S^2\cdot s_{\mathrm{dtype}}.
\]

例如：

\[
B=8,\quad n_h=32,\quad S=4096,\quad s_{\mathrm{dtype}}=2.
\]

则仅一个 $B\times n_h\times S\times S$ 矩阵就需要：

\[
8\times 32\times 4096^2\times 2
\approx
8.59\ \mathrm{GB}.
\]

这只是单层 attention 的一个中间矩阵。训练时还要考虑 backward 保存和重读，开销更加不可接受。

\begin{figure}[H]
\centering
\begin{tikzpicture}[
font=\small,
tensor/.style={draw, rounded corners=3pt, minimum width=1.55cm, minimum height=0.75cm, align=center},
big/.style={draw, rounded corners=3pt, minimum width=2.0cm, minimum height=1.15cm, align=center},
arrow/.style={-{Latex[length=2.2mm]}, thick}
]

\node[font=\bfseries] at (0,3.35) {标准 Attention 显式物化 $S\times S$ 中间矩阵};

\node[tensor, fill=blue!8] (q) at (-5.2,1.4) {$Q$\\$S,d_h$};
\node[tensor, fill=blue!8] (k) at (-5.2,0.35) {$K$\\$S,d_h$};
\node[big, fill=orange!14] (score) at (-2.6,0.85) {$S=QK^T$\\$S,S$};
\node[big, fill=green!12] (prob) at (0.25,0.85) {$P=\softmax(S)$\\$S,S$};
\node[tensor, fill=blue!8] (v) at (2.65,0.2) {$V$\\$S,d_h$};
\node[tensor, fill=purple!10] (o) at (5.0,0.85) {$O=PV$\\$S,d_h$};

\draw[arrow] (q) -- (score);
\draw[arrow] (k) -- (score);
\draw[arrow] (score) -- node[above] {write/read HBM} (prob);
\draw[arrow] (prob) -- (o);
\draw[arrow] (v) -- (o);

\node[draw, rounded corners=3pt, fill=red!8, minimum width=2.6cm, minimum height=0.55cm, align=center] at (-1.2,-0.95)
  {large $S\times S$ traffic};

\node[align=left, text width=10.5cm] at (0,-2.05) {
  标准实现往往把 score matrix 和 softmax matrix 写入 HBM，再从 HBM 读回。
  当序列长度变大时，中间矩阵读写成为主要瓶颈。
};

\end{tikzpicture}
\caption{标准 Attention 的中间矩阵物化流程}
\label{fig:standard-attention-materialization}
\end{figure}

\subsection{HBM 与 SRAM 的带宽差异}
\label{subsec:hbm-vs-sram-bandwidth}

现代 GPU 有多级存储层次。对深度学习 kernel 来说，最重要的几层是：

\begin{itemize}
\item \textbf{HBM}：容量大，带宽高，但相比片上存储仍然慢；
\item \textbf{L2 Cache}：片上共享缓存，容量比 HBM 小得多；
\item \textbf{Shared Memory / SRAM}：每个 SM 内部可编程片上存储，延迟低、带宽高、容量有限；
\item \textbf{Registers}：线程私有，最快但数量有限。
\end{itemize}

标准 Attention 的问题是，它在 HBM 中读写了太多中间结果。即使 GPU 的 Tensor Core 对矩阵乘法很快，如果数据需要反复从 HBM 搬进搬出，整体性能仍会被内存带宽限制。

可以把一个 GPU kernel 的执行理解为：

\[
\text{HBM}
\rightarrow
\text{SRAM/register}
\rightarrow
\text{compute}
\rightarrow
\text{HBM}.
\]

如果中间矩阵 $\mat{S}$ 和 $\mat{P}$ 都写回 HBM，再读回继续计算，则数据路径变成：

\[
Q,K,V
\rightarrow
\mat{S}\ \text{write HBM}
\rightarrow
\mat{S}\ \text{read HBM}
\rightarrow
\mat{P}\ \text{write HBM}
\rightarrow
\mat{P}\ \text{read HBM}
\rightarrow
O.
\]

FlashAttention 的目标是减少这类 HBM 往返。它将 Q、K、V 按块搬入 SRAM/register，在片上计算 score、softmax 和局部输出累积，中间 attention matrix 不落地到 HBM。

\begin{figure}[H]
\centering
\begin{tikzpicture}[
font=\small,
mem/.style={draw, rounded corners=4pt, minimum width=3.0cm, minimum height=0.75cm, align=center},
arrow/.style={-{Latex[length=2.2mm]}, thick}
]

\node[font=\bfseries] at (0,3.25) {GPU 存储层次与 Attention IO};

\node[mem, fill=orange!14] (hbm) at (0,2.0) {HBM\\large capacity, high latency};
\node[mem, fill=green!12] (l2) at (0,0.95) {L2 Cache};
\node[mem, fill=blue!8] (smem) at (0,-0.1) {Shared Memory / SRAM\\small capacity, high bandwidth};
\node[mem, fill=purple!10] (reg) at (0,-1.15) {Registers\\fastest, tiny};

\draw[arrow] (hbm) -- (l2);
\draw[arrow] (l2) -- (smem);
\draw[arrow] (smem) -- (reg);

\node[align=left, text width=10.5cm] at (0,-2.35) {
  FlashAttention 的关键是把尽可能多的中间计算留在 SRAM/register 中完成，
  避免把 $S\times S$ attention matrix 写入 HBM。
};

\end{tikzpicture}
\caption{HBM 与片上 SRAM/register 的存储层次}
\label{fig:hbm-sram-memory-hierarchy}
\end{figure}

\subsection{计算复杂度与 IO 复杂度}
\label{subsec:compute-complexity-and-io-complexity}

标准 Attention 和 FlashAttention 在数学上都计算精确 Attention。它们的 FLOPs 复杂度同阶：

\[
O(S^2d_h).
\]

这意味着 FlashAttention 并不是通过把 attention 变成稀疏或近似来减少数学计算。它的收益来自 IO 复杂度下降。

为了理解这一点，需要区分：

\begin{itemize}
\item \textbf{计算复杂度}：执行多少乘加、指数、除法等运算；
\item \textbf{IO 复杂度}：数据在 HBM 与片上存储之间搬运多少次。
\end{itemize}

在 GPU 上，如果 kernel 是 compute-bound，优化 FLOPs 更重要；如果 kernel 是 memory-bound，减少 HBM 读写更重要。标准 Attention 中，$S\times S$ 中间矩阵导致 HBM traffic 很高，尤其在长序列下，IO 成本可能主导运行时间。

粗略看，标准 Attention 至少涉及：

\[
\text{read }Q,K,V,
\]
\[
\text{write }S,
\]
\[
\text{read }S,
\]
\[
\text{write }P,
\]
\[
\text{read }P,V,
\]
\[
\text{write }O.
\]

其中 $S$ 和 $P$ 都是 $S\times S$ 大矩阵。FlashAttention 则尽量只读 Q/K/V，写 O，并保存少量 softmax 统计量，例如每行最大值和归一化分母：

\[
m_i,\quad l_i.
\]

因此，它把 attention 的显存复杂度从需要保存 $O(S^2)$ 中间矩阵，降低到只需要 $O(S)$ 级别的辅助统计量和 $O(Sd_h)$ 的输出。

\begin{table}[H]
\centering
\small
\caption{标准 Attention 与 FlashAttention 的关键区别}
\label{tab:standard-vs-flash-attention}
\begin{tabular}{p{0.24\textwidth}p{0.32\textwidth}p{0.34\textwidth}}
\toprule
\textbf{维度} & \textbf{标准 Attention} & \textbf{FlashAttention} \
\midrule
数学结果
& 精确 Attention
& 精确 Attention \
\midrule
FLOPs 复杂度
& $O(S^2d_h)$
& $O(S^2d_h)$，可能有少量重计算 \
\midrule
中间矩阵
& 常物化 $S\times S$ score/probability
& 不物化完整 attention matrix \
\midrule
主要优化目标
& 简单分步执行
& 减少 HBM 读写，提升 IO efficiency \
\midrule
显存占用
& 需要保存较大中间矩阵
& 只保存输出和少量 softmax 统计量 \
\midrule
适用价值
& 短序列或简单实现
& 长序列、高 batch、训练和 prefill 场景 \
\bottomrule
\end{tabular}
\end{table}

\subsection{为什么不是 FLOPs 越少越快}
\label{subsec:why-fewer-flops-not-always-faster}

很多性能优化初学者会自然地认为：算法 FLOPs 越少，速度越快。但在 GPU 上，这并不总成立。一个 kernel 的速度常常受以下因素共同决定：

\begin{itemize}
\item HBM 读写量；
\item L2 cache 命中；
\item shared memory 使用；
\item register pressure；
\item Tensor Core 利用率；
\item warp occupancy；
\item memory coalescing；
\item kernel launch 和同步；
\item 数据 layout；
\item 是否有中间 tensor 落地。
\end{itemize}

FlashAttention 的一个重要启发是：\textbf{宁可做一些额外计算，也要避免昂贵的 HBM 读写}。尤其在 backward 中，FlashAttention 可以选择不保存完整 attention probability，而是在 backward 时重新计算需要的 score 和 softmax 块。这增加了 FLOPs，但减少了巨大的中间矩阵存储和读取。

这与 gradient checkpointing 的思想类似：用重计算换显存和 IO。区别在于 FlashAttention 的重计算发生在更细粒度的 attention kernel 内部，并且与 tiling 和 online softmax 结合。

\begin{figure}[H]
\centering
\begin{tikzpicture}[
font=\small,
axis/.style={-{Latex[length=2.2mm]}, thick},
point/.style={circle, draw, fill=blue!10, minimum size=0.18cm},
line/.style={thick}
]

\node[font=\bfseries] at (0,3.0) {Roofline 视角：Attention 可能受 IO 而非 FLOPs 限制};

\draw[axis] (-5.0,-1.2) -- (5.1,-1.2) node[right] {arithmetic intensity};
\draw[axis] (-5.0,-1.2) -- (-5.0,2.2) node[above] {performance};

\draw[line, orange!80!black] (-4.7,-1.0) -- (-1.5,1.5);
\draw[line, orange!80!black] (-1.5,1.5) -- (4.7,1.5);

\node[point, fill=red!18] (std) at (-3.2,0.15) {};
\node[anchor=west] at (-3.0,0.15) {standard attention};

\node[point, fill=green!18] (flash) at (0.1,1.2) {};
\node[anchor=west] at (0.3,1.2) {FlashAttention};

\node[align=center, font=\scriptsize] at (-3.2,-0.45) {memory-bound};
\node[align=center, font=\scriptsize] at (2.2,1.85) {compute-bound ceiling};

\node[align=left, text width=10.5cm] at (0,-2.35) {
  减少 HBM traffic 可以提高 arithmetic intensity，使 kernel 从 memory-bound 区域向更高性能区域移动。
  因此即使 FLOPs 没有减少，甚至略有增加，wall-clock 时间仍可能下降。
};

\end{tikzpicture}
\caption{从 roofline 模型理解 FlashAttention 的 IO-aware 思路}
\label{fig:attention-roofline-intuition}
\end{figure}

\section{Online Softmax}
\label{sec:online-softmax}

FlashAttention 的数学核心之一是 online softmax。Attention 的困难在于 softmax 的分母依赖整行所有元素：

\[
\softmax(x_i)
\frac{e^{x_i}}
{\sum_j e^{x_j}}.
\]

如果把 score matrix 按列分块计算，那么在只看到一个块时，并不知道整行的全局最大值和全局分母。Online softmax 解决了这个问题：它允许我们一块一块处理 score，同时维护每一行的最大值和指数和，最终得到与一次性 softmax 完全一致的结果。

\subsection{softmax 的数值稳定性}
\label{subsec:softmax-numerical-stability}

直接计算：

\[
\frac{e^{x_i}}{\sum_j e^{x_j}}
\]

容易溢出。数值稳定做法是减去最大值：

\[
m=\max_j x_j.
\]

然后：

\[
\softmax(x_i)
\frac{e^{x_i-m}}
{\sum_j e^{x_j-m}}.
\]

因为：

\[
\frac{e^{x_i-m}}{\sum_j e^{x_j-m}}
# \frac{e^{x_i}e^{-m}}{\sum_j e^{x_j}e^{-m}}
\frac{e^{x_i}}{\sum_j e^{x_j}}.
\]

减去 $m$ 不改变 softmax 结果，但能防止指数溢出。

标准 softmax 通常分两步：

\[
m=\max_j x_j,
\]
\[
l=\sum_j e^{x_j-m}.
\]

然后输出：

\[
y_i=\frac{e^{x_i-m}}{l}.
\]

如果 $x$ 是 attention score 的一行，那么 $j$ 遍历所有 key positions。FlashAttention 分块处理时，需要在多个 blocks 之间维护同样的 $m$ 和 $l$。

\subsection{分块计算中的 max 与 sum 更新}
\label{subsec:blockwise-max-sum-update}

假设一行 score 被分成多个 blocks。当前已经处理过若干 blocks，维护：

\[
m_{\mathrm{old}}
\max_{\text{old blocks}} x_j,
\]
\[
l_{\mathrm{old}}
\sum_{\text{old blocks}} e^{x_j-m_{\mathrm{old}}}.
\]

现在处理新 block，计算该 block 的局部最大值：

\[
m_{\mathrm{blk}}
\max_{\text{new block}} x_j,
\]

和局部指数和：

\[
l_{\mathrm{blk}}
\sum_{\text{new block}} e^{x_j-m_{\mathrm{blk}}}.
\]

合并 old blocks 与 new block 的全局最大值为：

\[
m_{\mathrm{new}}
\max(m_{\mathrm{old}},m_{\mathrm{blk}}).
\]

旧的指数和 $l_{\mathrm{old}}$ 是以 $m_{\mathrm{old}}$ 为基准的。要改成以 $m_{\mathrm{new}}$ 为基准，需要乘缩放因子：

\[
e^{m_{\mathrm{old}}-m_{\mathrm{new}}}.
\]

新 block 的指数和也要改成以 $m_{\mathrm{new}}$ 为基准：

\[
e^{m_{\mathrm{blk}}-m_{\mathrm{new}}}.
\]

因此：

\[
l_{\mathrm{new}}
e^{m_{\mathrm{old}}-m_{\mathrm{new}}}l_{\mathrm{old}}
+
e^{m_{\mathrm{blk}}-m_{\mathrm{new}}}l_{\mathrm{blk}}.
\]

这就是 online softmax 的核心更新公式。

\subsection{避免完整 attention matrix}
\label{subsec:avoid-full-attention-matrix}

Attention 输出不是 softmax 本身，而是：

\[
\vect{o}
\sum_j
\frac{e^{x_j-m}}{l}
\vect{v}_j.
\]

分块时，不仅要维护 $m$ 和 $l$，还要维护未归一化的输出累积。设旧输出累积为：

\[
\vect{o}_{\mathrm{old}}
\sum_{\text{old blocks}}
\frac{e^{x_j-m_{\mathrm{old}}}}
{l_{\mathrm{old}}}
\vect{v}_j.
\]

更方便的做法是维护带分母的累积量：

\[
\vect{a}_{\mathrm{old}}
# l_{\mathrm{old}}\vect{o}_{\mathrm{old}}
\sum_{\text{old blocks}}
e^{x_j-m_{\mathrm{old}}}\vect{v}_j.
\]

新 block 的累积量为：

\[
\vect{a}_{\mathrm{blk}}
\sum_{\text{new block}}
e^{x_j-m_{\mathrm{blk}}}\vect{v}_j.
\]

合并到新基准 $m_{\mathrm{new}}$ 后：

\[
\vect{a}_{\mathrm{new}}
e^{m_{\mathrm{old}}-m_{\mathrm{new}}}\vect{a}*{\mathrm{old}}
+
e^{m*{\mathrm{blk}}-m_{\mathrm{new}}}\vect{a}_{\mathrm{blk}}.
\]

最终输出：

\[
\vect{o}_{\mathrm{new}}
\frac{\vect{a}*{\mathrm{new}}}{l*{\mathrm{new}}}.
\]

这意味着我们可以一块一块处理 score 和 V，持续更新：

\[
m_i,\quad l_i,\quad \vect{a}_i,
\]

而不需要保存整行 $S$ 个 softmax probability，更不需要保存完整 $S\times S$ attention matrix。

\begin{figure}[H]
\centering
\begin{tikzpicture}[
font=\small,
block/.style={draw, rounded corners=3pt, minimum width=1.45cm, minimum height=0.65cm, align=center},
state/.style={draw, rounded corners=3pt, minimum width=1.4cm, minimum height=0.55cm, align=center, fill=green!10},
arrow/.style={-{Latex[length=2.1mm]}, thick}
]

\node[font=\bfseries] at (0,3.2) {Online Softmax 分块更新};

\node[block, fill=blue!8] (b0) at (-5.1,1.3) {score\\block 0};
\node[state] (s0) at (-3.3,1.3) {$m,l,a$};

\node[block, fill=blue!8] (b1) at (-1.7,1.3) {score\\block 1};
\node[state] (s1) at (0.1,1.3) {$m,l,a$};

\node[block, fill=blue!8] (b2) at (1.7,1.3) {score\\block 2};
\node[state] (s2) at (3.5,1.3) {$m,l,a$};

\node[block, fill=purple!10] (out) at (5.2,1.3) {$O=a/l$};

\draw[arrow] (b0) -- (s0);
\draw[arrow] (s0) -- (b1);
\draw[arrow] (b1) -- (s1);
\draw[arrow] (s1) -- (b2);
\draw[arrow] (b2) -- (s2);
\draw[arrow] (s2) -- (out);

\node[align=left, text width=10.5cm] at (0,-0.65) {
  每处理一个 score block，就更新该行的最大值 $m$、指数和 $l$ 与输出累积 $a$。
  最终结果与一次性 softmax 完全一致，但不需要物化完整 attention row。
};

\end{tikzpicture}
\caption{Online softmax 的分块归约流程}
\label{fig:online-softmax-blockwise}
\end{figure}

\subsection{online softmax 推导}
\label{subsec:online-softmax-derivation}

下面把推导整理成更系统的形式。对某一行 attention score，分块处理。第 $r$ 个 block 的元素集合为 $\mathcal{B}_r$。

局部最大值：

\[
m_r=\max_{j\in\mathcal{B}_r} x_j.
\]

局部指数和：

\[
l_r=\sum_{j\in\mathcal{B}_r}e^{x_j-m_r}.
\]

局部 value 累积：

\[
\vect{a}_r
\sum_{j\in\mathcal{B}_r}
e^{x_j-m_r}\vect{v}_j.
\]

若当前已经合并到第 $r-1$ 个 block，状态为：

\[
m^{(r-1)},\quad l^{(r-1)},\quad \vect{a}^{(r-1)}.
\]

合并第 $r$ 个 block：

\[
m^{(r)}
\max(m^{(r-1)},m_r).
\]

\[
l^{(r)}
e^{m^{(r-1)}-m^{(r)}}l^{(r-1)}
+
e^{m_r-m^{(r)}}l_r.
\]

\[
\vect{a}^{(r)}
e^{m^{(r-1)}-m^{(r)}}\vect{a}^{(r-1)}
+
e^{m_r-m^{(r)}}\vect{a}_r.
\]

处理完所有 blocks 后：

\[
\vect{o}
\frac{\vect{a}^{(R)}}{l^{(R)}}.
\]

这个公式的关键在于缩放因子。每当全局最大值变大，旧累积量必须被重新缩放，否则 softmax 分母不一致。FlashAttention 的 kernel 就是把这种数学更新嵌入到 tiled attention 计算中。

下面给出一个教学版 online softmax attention 的 Python 伪代码。它强调数学流程，不代表高性能实现。

\begin{lstlisting}[language=Python,caption={教学版 Online Softmax Attention 伪代码},label={lst:online-softmax-attention-pseudocode}]
import math
import torch

def online_attention_one_row(q, k, v, block_size):
"""
q: [d]
k: [seq_len, d]
v: [seq_len, d]
return: [d]
"""
seq_len = k.shape[0]
scale = 1.0 / math.sqrt(q.shape[0])

m = torch.tensor(-float("inf"), device=q.device)
l = torch.tensor(0.0, device=q.device)
acc = torch.zeros_like(q)

for start in range(0, seq_len, block_size):
    end = min(start + block_size, seq_len)

    k_blk = k[start:end]  # [B_blk, d]
    v_blk = v[start:end]  # [B_blk, d]

    scores = (k_blk @ q) * scale  # [B_blk]
    m_blk = torch.max(scores)
    p_blk = torch.exp(scores - m_blk)
    l_blk = torch.sum(p_blk)
    acc_blk = p_blk @ v_blk

    m_new = torch.maximum(m, m_blk)

    old_scale = torch.exp(m - m_new)
    blk_scale = torch.exp(m_blk - m_new)

    acc = old_scale * acc + blk_scale * acc_blk
    l = old_scale * l + blk_scale * l_blk
    m = m_new

return acc / l

\end{lstlisting}

实际 FlashAttention 会同时处理一个 $Q$ block 的多行，而不是一行；会在 shared memory/register 中计算 block GEMM；会进行向量化、warp-level reduction 和 Tensor Core MMA。上面代码只是展示 online softmax 的数学骨架。

\section{FlashAttention 核心算法}
\label{sec:flashattention-core-algorithm}

FlashAttention 的核心可以概括为四句话：

\begin{enumerate}
\item 将 $Q,K,V$ 沿序列维分块；
\item 将 $K,V$ blocks 搬入 SRAM；
\item 对 $Q$ blocks 计算局部 attention，并用 online softmax 更新输出；
\item 不把完整 $S\times S$ attention matrix 写入 HBM。
\end{enumerate}

它的设计目标不是减少 Attention 的数学复杂度，而是最小化 HBM 与片上存储之间的 IO。

\subsection{tiling}
\label{subsec:flashattention-tiling}

设序列长度为 $S$。将 Q 按行切成 blocks：

\[
\mat{Q}_1,\mat{Q}*2,\ldots,\mat{Q}*{T_Q},
\]

每个 Q block 大小约为：

\[
B_Q\times d_h.
\]

将 K/V 也按序列维切成 blocks：

\[
(\mat{K}_1,\mat{V}_1),
(\mat{K}*2,\mat{V}*2),
\ldots,
(\mat{K}*{T_K},\mat{V}*{T_K}),
\]

每个 K/V block 大小约为：

\[
B_K\times d_h.
\]

对于一个 Q block 和一个 K block，局部 score 为：

\[
\mat{S}_{ij}
\mat{Q}_i\mat{K}_j^{\mathsf{T}}.
\]

它的 shape 为：

\[
B_Q\times B_K.
\]

这个局部矩阵可以放在 SRAM/register 中处理，而不需要写入 HBM。对所有 K/V blocks 迭代，使用 online softmax 更新 Q block 对应的输出。

\[
\mat{O}_i
\operatorname{Attention}(\mat{Q}_i,\mat{K},\mat{V}).
\]

最终只将 $\mat{O}_i$ 写回 HBM。

\begin{figure}[H]
\centering
\begin{tikzpicture}[
font=\small,
tile/.style={draw, rounded corners=2pt, minimum width=0.9cm, minimum height=0.55cm, align=center},
mat/.style={draw, rounded corners=3pt, minimum width=1.5cm, minimum height=1.0cm, align=center},
arrow/.style={-{Latex[length=2.1mm]}, thick}
]

\node[font=\bfseries] at (0,3.35) {FlashAttention Tiling 数据流};

\node[mat, fill=blue!8] (q) at (-5.2,1.25) {$Q$};
\foreach \i/\y in {0/1.75,1/1.25,2/0.75} {
  \node[tile, fill=blue!12] at (-5.2,\y) {$Q_{\i}$};
}

\node[mat, fill=green!10] (kv) at (-2.8,1.25) {$K,V$};
\foreach \i/\y in {0/1.75,1/1.25,2/0.75} {
  \node[tile, fill=green!14] at (-2.8,\y) {$K_{\i},V_{\i}$};
}

\node[mat, fill=orange!14, minimum width=2.2cm] (sram) at (0.0,1.25) {SRAM/register\\block compute};

\node[mat, fill=purple!10] (online) at (2.9,1.25) {online\\softmax update};

\node[mat, fill=red!8] (out) at (5.2,1.25) {$O$\\write once};

\draw[arrow] (q) -- (sram);
\draw[arrow] (kv) -- (sram);
\draw[arrow] (sram) -- (online);
\draw[arrow] (online) -- (out);

\node[align=left, text width=10.5cm] at (0,-0.85) {
  $Q,K,V$ 按块搬入片上存储。局部 score、mask、softmax 和 $PV$ 累积都在片上完成。
  完整 $S\times S$ attention matrix 不写入 HBM。
};

\end{tikzpicture}
\caption{FlashAttention 的 tiling 数据流}
\label{fig:flashattention-tiling-dataflow}
\end{figure}

\subsection{SRAM 中计算 block}
\label{subsec:compute-block-in-sram}

FlashAttention 的一个典型循环结构可以抽象为：

\begin{enumerate}
\item 外层遍历 K/V blocks；
\item 内层遍历 Q blocks；
\item 将当前 K/V block 加载到 SRAM；
\item 将当前 Q block 加载到 SRAM/register；
\item 计算局部 score；
\item 应用 causal mask 或 attention mask；
\item 计算局部 softmax 统计；
\item 更新输出累积；
\item 最后写回 O。
\end{enumerate}

不同实现可能选择不同循环顺序。关键原则是：让 K/V block 被加载后尽可能复用，减少 HBM 访问；让局部 score matrix 只在片上存在。

对一个 Q block，维护每一行的：

\[
m_i,\quad l_i,\quad \vect{a}_i.
\]

每处理一个 K/V block，更新这些状态。最后输出：

\[
\vect{o}_i=\frac{\vect{a}_i}{l_i}.
\]

如果是 causal attention，当 K block 的位置超过 Q block 当前行允许关注的位置，需要 mask 掉未来位置：

\[
S_{ij}=-\infty,\quad j>i.
\]

在 block 级别，causal mask 可能出现三种情况：

\begin{itemize}
\item 当前 K block 完全在 Q block 之前：无需 mask；
\item 当前 K block 完全在 Q block 之后：整块跳过；
\item 当前 K block 与 Q block 对角线相交：需要块内 mask。
\end{itemize}

这种 block-level 判断可以避免不必要计算。

\subsection{fused kernel}
\label{subsec:flashattention-fused-kernel}

标准实现可能将 Attention 拆成多个 kernel：

\[
QK^{\mathsf{T}}\ \text{kernel},
\]
\[
mask/scale\ \text{kernel},
\]
\[
softmax\ \text{kernel},
\]
\[
dropout\ \text{kernel},
\]
\[
PV\ \text{kernel}.
\]

每个 kernel 都可能把中间结果写回 HBM，再被下一个 kernel 读入。FlashAttention 则将这些操作融合到一个或少数几个 kernel 中：

\[
QK^{\mathsf{T}}
+
mask
+
softmax
+
dropout
+
PV
\rightarrow
\text{fused attention kernel}.
\]

Kernel fusion 的收益包括：

\begin{itemize}
\item 减少 kernel launch overhead；
\item 减少 HBM 中间读写；
\item 中间值停留在 register/shared memory；
\item 更容易做跨操作优化；
\item 更低显存峰值。
\end{itemize}

但 fused kernel 也更复杂。它需要同时管理矩阵乘法、softmax、mask、dropout、数值稳定、寄存器压力、shared memory、warp 分工和边界条件。FlashAttention 的难点正是把这些操作合并后仍然保持高性能。

\begin{figure}[H]
\centering
\begin{tikzpicture}[
font=\small,
box/.style={draw, rounded corners=3pt, minimum width=1.45cm, minimum height=0.58cm, align=center},
arrow/.style={-{Latex[length=2.1mm]}, thick}
]

\node[font=\bfseries] at (0,3.1) {Kernel Fusion：从多次 HBM 往返到单个 fused kernel};

% standard
\node[font=\bfseries] at (-3.5,2.35) {Standard};
\node[box, fill=blue!8] (a1) at (-5.5,1.45) {$QK^T$};
\node[box, fill=orange!14] (a2) at (-3.8,1.45) {mask};
\node[box, fill=green!10] (a3) at (-2.1,1.45) {softmax};
\node[box, fill=purple!10] (a4) at (-0.4,1.45) {$PV$};
\draw[arrow] (a1) -- node[above,font=\scriptsize] {HBM} (a2);
\draw[arrow] (a2) -- node[above,font=\scriptsize] {HBM} (a3);
\draw[arrow] (a3) -- node[above,font=\scriptsize] {HBM} (a4);

% flash
\node[font=\bfseries] at (2.5,2.35) {FlashAttention};
\node[box, fill=red!8, minimum width=3.0cm, minimum height=0.85cm] (fused) at (3.2,1.45)
  {$QK^T$ + mask + softmax + $PV$};
\node[box, fill=yellow!16] (o) at (5.7,1.45) {$O$};
\draw[arrow] (fused) -- node[above,font=\scriptsize] {write once} (o);

\node[align=left, text width=10.5cm] at (0,-0.45) {
  标准实现的多个 kernel 之间需要通过 HBM 传递中间矩阵。
  FlashAttention 把这些步骤融合，让中间 attention block 不落地。
};

\end{tikzpicture}
\caption{FlashAttention 的 fused kernel 思路}
\label{fig:flashattention-fused-kernel}
\end{figure}

\subsection{forward pass}
\label{subsec:flashattention-forward-pass}

FlashAttention forward 的核心状态包括：

\[
\mat{O}\in\mathbb{R}^{S\times d_h},
\]
\[
\vect{m}\in\mathbb{R}^{S},
\]
\[
\vect{l}\in\mathbb{R}^{S}.
\]

其中 $\vect{m}$ 保存每一行的运行最大值，$\vect{l}$ 保存每一行的运行 softmax 分母。输出 $\mat{O}$ 在分块过程中逐步更新。

抽象伪代码如下：

\begin{lstlisting}[language=Python,caption={FlashAttention forward 的简化伪代码},label={lst:flashattention-forward-pseudocode}]
def flash_attention_forward(Q, K, V, block_q, block_k):
"""
Q, K, V: [seq_len, head_dim]
This is a mathematical sketch, not an efficient implementation.
"""
S, D = Q.shape

O = zeros([S, D])
m = full([S], -inf)
l = zeros([S])

for q_start in range(0, S, block_q):
    q_end = min(q_start + block_q, S)
    Q_blk = Q[q_start:q_end]

    O_blk = zeros([q_end - q_start, D])
    m_blk_running = full([q_end - q_start], -inf)
    l_blk_running = zeros([q_end - q_start])

    for k_start in range(0, S, block_k):
        k_end = min(k_start + block_k, S)
        K_blk = K[k_start:k_end]
        V_blk = V[k_start:k_end]

        scores = Q_blk @ K_blk.T / sqrt(D)

        # Apply causal mask if needed.
        scores = apply_causal_mask(
            scores,
            q_start=q_start,
            k_start=k_start,
        )

        m_local = row_max(scores)
        p = exp(scores - m_local[:, None])
        l_local = row_sum(p)
        acc_local = p @ V_blk

        m_new = maximum(m_blk_running, m_local)
        old_scale = exp(m_blk_running - m_new)
        new_scale = exp(m_local - m_new)

        O_blk = (
            old_scale[:, None] * O_blk
            + new_scale[:, None] * acc_local
        )
        l_blk_running = (
            old_scale * l_blk_running
            + new_scale * l_local
        )
        m_blk_running = m_new

    O_blk = O_blk / l_blk_running[:, None]

    O[q_start:q_end] = O_blk
    m[q_start:q_end] = m_blk_running
    l[q_start:q_end] = l_blk_running

return O, m, l

\end{lstlisting}

实际 kernel 不会像 Python 这样分配中间张量。它会把 block 内矩阵乘法映射到 GPU thread block、warp 和 Tensor Core MMA 上，尽量保持 Q/K/V block 的加载、计算和累积在片上完成。

\subsection{backward pass}
\label{subsec:flashattention-backward-pass}

训练时还需要 backward。标准 Attention backward 如果保存了 $\mat{P}$，可以直接使用它计算梯度。但 $\mat{P}$ 是 $S\times S$ 大矩阵，保存它会带来高显存和 IO。

FlashAttention 的 backward 通常不保存完整 $\mat{P}$，而是保存 forward 中每行的 softmax 统计量：

\[
m_i,\quad l_i.
\]

在 backward 中，对每个 block 重新计算：

\[
\mat{S}_{ij}=\mat{Q}_i\mat{K}*j^{\mathsf{T}},
\]
\[
\mat{P}*{ij}
\frac{\exp(\mat{S}_{ij}-m_i)}{l_i}.
\]

然后使用重算得到的 $\mat{P}_{ij}$ 计算：

\[
d\mat{V},
\quad
d\mat{Q},
\quad
d\mat{K}.
\]

这体现了重计算换 IO 的思想。Backward 的 FLOPs 增加了一些，但避免读取和保存巨大的 attention probability matrix。

Softmax backward 的核心公式为：

\[
dS_i
P_i
\odot
\left(
dP_i
-
\sum_j dP_{ij}P_{ij}
\right).
\]

在 Attention 中：

\[
O=PV.
\]

因此：

\[
dP=dO V^{\mathsf{T}},
\]
\[
dV=P^{\mathsf{T}}dO.
\]

再由：

\[
S=QK^{\mathsf{T}},
\]

得到：

\[
dQ=dS K,
\]
\[
dK=dS^{\mathsf{T}}Q.
\]

FlashAttention backward 通过分块重算 $P$，在片上完成这些局部梯度计算，并将 $dQ,dK,dV$ 写回或累积。

\begin{figure}[H]
\centering
\begin{tikzpicture}[
font=\small,
box/.style={draw, rounded corners=3pt, minimum width=1.9cm, minimum height=0.7cm, align=center},
arrow/.style={-{Latex[length=2.1mm]}, thick}
]

\node[font=\bfseries] at (0,3.25) {FlashAttention Backward：重算 Attention Block 而非保存完整 $P$};

\node[box, fill=blue!8] (saved) at (-5.0,1.4) {saved\\$m,l$};
\node[box, fill=green!10] (qkv) at (-5.0,0.35) {$Q,K,V$\\blocks};
\node[box, fill=orange!14] (recompute) at (-2.3,0.9) {recompute\\$S,P$ block};
\node[box, fill=purple!10] (grad) at (0.5,0.9) {local gradient\\$dQ,dK,dV$};
\node[box, fill=red!8] (acc) at (3.3,0.9) {accumulate\\write grads};

\draw[arrow] (saved) -- (recompute);
\draw[arrow] (qkv) -- (recompute);
\draw[arrow] (recompute) -- (grad);
\draw[arrow] (grad) -- (acc);

\node[align=left, text width=10.5cm] at (0,-1.2) {
  Forward 只保存每行 softmax 统计量 $m,l$。Backward 中分块重算 $S$ 和 $P$，
  用额外计算换取更低显存与更少 HBM 读取。
};

\end{tikzpicture}
\caption{FlashAttention backward 的重计算思想}
\label{fig:flashattention-backward-recompute}
\end{figure}

\subsection{FlashAttention-2 / FlashAttention-3 的优化方向}
\label{subsec:flashattention2-flashattention3}

FlashAttention 的第一版解决了 IO-aware exact attention 的核心问题，但仍然距离高度优化的 GEMM 有差距。后续版本主要围绕 GPU 并行度、work partitioning、硬件特性和低精度展开。

\textbf{FlashAttention-2} 的核心方向包括：

\begin{itemize}
\item 减少非矩阵乘法 FLOPs；
\item 改进 thread block 和 warp 之间的工作划分；
\item 对单个 head 的 attention 进行更好的并行切分；
\item 降低 shared memory 读写和跨 warp 通信；
\item 提高 occupancy 和 Tensor Core 利用率。
\end{itemize}

从系统角度看，FlashAttention-2 的启发是：减少 IO 只是第一步。一个 kernel 想接近 GEMM 性能，还必须让 GPU 有足够并行度，避免某些 warp 空闲，避免 shared memory 成为内部瓶颈。

\textbf{FlashAttention-3} 则进一步面向新硬件特性，尤其是 Hopper 架构中的异步数据搬运、Tensor Memory Accelerator、warp specialization 和 FP8 支持。其思想可以概括为：

\begin{itemize}
\item 让数据搬运和计算重叠；
\item 让不同 warp 承担不同角色，例如 producer/consumer；
\item 更细粒度地交错矩阵乘法和 softmax；
\item 利用 FP8 等低精度能力提升吞吐；
\item 在低精度下控制数值误差。
\end{itemize}

对 AI Infra 工程师来说，理解这些方向比记住某个版本的具体实现更重要。FlashAttention 的演进说明，高性能 attention kernel 不是单一算法公式，而是算法、数值稳定、GPU 存储层次、线程组织、低精度和硬件特性的共同设计。

\begin{table}[H]
\centering
\small
\caption{FlashAttention 系列的优化演进}
\label{tab:flashattention-evolution}
\begin{tabular}{p{0.22\textwidth}p{0.34\textwidth}p{0.34\textwidth}}
\toprule
\textbf{版本} & \textbf{核心问题} & \textbf{主要优化方向} \
\midrule
FlashAttention
& 标准 Attention 物化 $S\times S$ 中间矩阵，HBM IO 过高
& IO-aware tiling、online softmax、不落地 attention matrix。 \
\midrule
FlashAttention-2
& 第一版仍未充分接近 GEMM 效率
& 更好的 work partitioning、更高 occupancy、减少非 matmul FLOPs。 \
\midrule
FlashAttention-3
& 新硬件特性未充分利用
& Hopper 异步、warp specialization、TMA、FP8 与低精度数值控制。 \
\bottomrule
\end{tabular}
\end{table}

\section{工程实现视角}
\label{sec:flashattention-engineering-view}

FlashAttention 的论文公式展示了算法思想，但工程实现还要处理大量底层细节。一个高性能 attention kernel 需要同时满足：

\begin{itemize}
\item 数据从 HBM 读取次数少；
\item shared memory 容量足够；
\item register pressure 不过高；
\item Tensor Core 利用率高；
\item warp 分工合理；
\item occupancy 足够；
\item 支持 causal mask、padding mask、dropout；
\item 支持不同 head dimension；
\item 支持 MHA/MQA/GQA；
\item 支持 variable sequence length；
\item 支持训练 backward 和推理 decode/prefill；
\item 与 PyTorch、Triton、CUDA graph、serving engine 集成。
\end{itemize}

本节从几个工程维度展开。

\subsection{block size 选择}
\label{subsec:flashattention-block-size-choice}

FlashAttention 的 block size 决定每次搬入 SRAM 的 Q/K/V tile 大小。设 Q block 大小为 $B_Q$，K/V block 大小为 $B_K$。局部 score block 为：

\[
B_Q\times B_K.
\]

Block size 受以下约束：

\begin{itemize}
\item shared memory 容量；
\item register 数量；
\item head dimension $d_h$；
\item Tensor Core tile 形状；
\item warp 数量；
\item occupancy；
\item causal mask 边界；
\item sequence length；
\item dtype。
\end{itemize}

如果 block 太大：

\begin{itemize}
\item shared memory 不够；
\item register pressure 高；
\item occupancy 下降；
\item block 内同步和 reduction 变重。
\end{itemize}

如果 block 太小：

\begin{itemize}
\item 数据复用不足；
\item kernel launch 和 block 调度开销相对变高；
\item Tensor Core 利用率可能下降；
\item online softmax 更新次数增加。
\end{itemize}

因此 block size 不是越大越好，也不是越小越好。高性能库通常会根据 head dimension、dtype、GPU 架构和是否 causal 选择不同 kernel specialization。

\subsection{warp-level parallelism}
\label{subsec:warp-level-parallelism}

在 GPU 中，warp 是调度基本单位。FlashAttention kernel 需要将一个 attention tile 的工作分配给多个 warps。例如：

\begin{itemize}
\item 一些 warps 负责加载 Q/K/V；
\item 一些 warps 负责 MMA；
\item 一些 warps 负责 softmax reduction；
\item 一些 warps 负责输出累积和写回。
\end{itemize}

如果 warp 分工不好，会出现：

\begin{itemize}
\item 某些 warp 等待其他 warp；
\item shared memory 读写冲突；
\item reduction 开销过大；
\item occupancy 低；
\item Tensor Core pipeline 填不满。
\end{itemize}

FlashAttention-2 的一个重要方向就是更好的 work partitioning。特别是在序列长度较短、batch/head 数不足或单 head 工作量较大时，需要在一个 head 内部进一步并行，否则 GPU 可能吃不满。

\begin{figure}[H]
\centering
\begin{tikzpicture}[
font=\small,
warp/.style={draw, rounded corners=3pt, minimum width=1.4cm, minimum height=0.55cm, align=center},
smem/.style={draw, rounded corners=4pt, minimum width=4.2cm, minimum height=0.75cm, align=center, fill=orange!14},
arrow/.style={-{Latex[length=2.0mm]}, thick}
]

\node[font=\bfseries] at (0,3.15) {FlashAttention kernel 中的 warp-level 分工示意};

\foreach \i/\x/\task/\col in {
  0/-4.5/load Q/blue!8,
  1/-2.6/load KV/green!10,
  2/-0.7/MMA/purple!10,
  3/1.2/softmax/yellow!16,
  4/3.1/accumulate/red!8
} {
  \node[warp, fill=\col] (w\i) at (\x,1.55) {Warp \i\\\task};
}

\node[smem] (smem) at (-0.7,0.25) {Shared Memory / Registers};

\foreach \i in {0,1,2,3,4} {
  \draw[arrow] (w\i) -- (smem);
}

\node[align=left, text width=10.5cm] at (0,-1.25) {
  实际实现中 warp 分工会更加复杂。核心目标是让数据搬运、矩阵乘法、
  softmax reduction 和输出累积尽量流水化，减少等待和片上通信。
};

\end{tikzpicture}
\caption{FlashAttention 中 warp-level parallelism 的工程直觉}
\label{fig:flashattention-warp-level-parallelism}
\end{figure}

\subsection{shared memory 使用}
\label{subsec:flashattention-shared-memory-usage}

Shared memory 是 FlashAttention 的关键资源。它用于缓存 K/V blocks、部分 Q blocks、中间 score 或 softmax 相关数据。使用 shared memory 的目标是减少 HBM 访问，但 shared memory 容量有限，且会影响 occupancy。

需要关注：

\begin{itemize}
\item 每个 thread block 使用多少 shared memory；
\item 是否导致每个 SM 可驻留 block 数下降；
\item shared memory bank conflict；
\item 数据 layout 是否利于 vectorized load；
\item Q/K/V tile 是否复用充分；
\item 是否需要 double buffering。
\end{itemize}

Double buffering 是常见优化：当当前 tile 正在计算时，异步加载下一个 tile。这样可以重叠数据搬运和计算。较新的 GPU 架构提供更强的异步拷贝能力，使这类优化更重要。

\subsection{与 Triton kernel 的关系}
\label{subsec:flashattention-and-triton}

Triton 是一种用于编写高性能 GPU kernel 的语言和编译器，常用于实现矩阵乘法、softmax、attention 等算子。FlashAttention 可以用 CUDA C++ 实现，也可以用 Triton 实现。Triton 的优势是：

\begin{itemize}
\item Python 风格开发体验；
\item 更容易写出 tile-based kernel；
\item 自动处理部分底层细节；
\item 与 PyTorch 集成方便；
\item 适合快速迭代和定制 kernel。
\end{itemize}

但极致性能实现仍可能需要 CUDA C++、CUTLASS、内联 PTX 或硬件特定优化。尤其当需要利用最新架构的异步拷贝、TMA、warp specialization、FP8 Tensor Core 等特性时，底层实现复杂度会增加。

对工程团队而言，选择 Triton 还是 CUDA C++ 取决于：

\begin{itemize}
\item 性能目标；
\item 开发速度；
\item 目标 GPU 架构；
\item 需要支持的 dtype 和 head dimension；
\item 是否需要快速适配模型变体；
\item 团队对底层 CUDA 的掌握程度。
\end{itemize}

下面给出一个简化的 Triton 风格伪代码，用于展示 tiled attention kernel 的形状。它不是可直接运行的完整实现。

\begin{lstlisting}[language=Python,caption={Triton 风格 FlashAttention kernel 的结构示意},label={lst:triton-flashattention-sketch}]

# This is a structural sketch, not runnable Triton code.

@triton_jit
def flash_attention_kernel(Q, K, V, O, LSE, stride_info, BLOCK_M, BLOCK_N):
# Program id selects a Q block and a head.
q_block_id = program_id(0)
head_id = program_id(1)

# Load Q tile.
q = load_q_tile(Q, q_block_id, head_id, BLOCK_M)

# Initialize online softmax states.
m = full([BLOCK_M], -inf)
l = zeros([BLOCK_M])
acc = zeros([BLOCK_M, HEAD_DIM])

# Loop over K/V blocks.
for kv_block_id in range(0, num_kv_blocks):
    k = load_k_tile(K, kv_block_id, head_id, BLOCK_N)
    v = load_v_tile(V, kv_block_id, head_id, BLOCK_N)

    scores = dot(q, transpose(k)) * scale
    scores = apply_mask_if_needed(scores, q_block_id, kv_block_id)

    m_new = maximum(m, max(scores, axis=1))
    p = exp(scores - m_new[:, None])

    alpha = exp(m - m_new)
    acc = acc * alpha[:, None] + dot(p, v)
    l = l * alpha + sum(p, axis=1)
    m = m_new

# Normalize and write output.
acc = acc / l[:, None]
store_o_tile(O, q_block_id, head_id, acc)

# Save log-sum-exp or m/l for backward if training.
store_lse(LSE, q_block_id, head_id, m + log(l))

\end{lstlisting}

这个伪代码保留了 FlashAttention 的核心结构：tile、循环、online softmax、累积输出、不保存 $S\times S$ 矩阵。

\section{FlashAttention 与推理系统}
\label{sec:flashattention-and-serving-systems}

FlashAttention 最初常被训练和长序列 attention 场景关注，但它同样影响推理系统。推理中有 prefill 和 decode 两种阶段，它们对 attention kernel 的需求不同。

\subsection{Prefill 中的 FlashAttention}
\label{subsec:flashattention-in-prefill}

Prefill 阶段一次处理 prompt 的全部 token。对于长 prompt，attention 仍然是 $S\times S$ 结构，因此 FlashAttention 非常有价值。它可以：

\begin{itemize}
\item 降低 prefill 显存峰值；
\item 减少 attention 中间矩阵 HBM 读写；
\item 提升长 prompt 的 TTFT；
\item 支持更长上下文；
\item 与 causal mask 结合。
\end{itemize}

在 serving 中，TTFT 通常由排队时间和 prefill 时间共同决定。长 prompt 请求尤其依赖高效 prefill attention。FlashAttention 可以显著降低 prefill kernel 的 IO 开销，但它不能消除 $O(S^2)$ 的数学计算，因此超长 prompt 仍然昂贵。对于超长上下文，还需要 chunked prefill、prefix cache、sparse attention 或 context parallel 等技术配合。

\subsection{Decode 中的 FlashAttention}
\label{subsec:flashattention-in-decode}

Decode 阶段每个请求通常只处理一个新 token。此时 Q 的 sequence length 为 $1$，K/V length 为历史上下文长度 $S$。Attention 形状是：

\[
Q: 1\times d_h,
\quad
K,V: S\times d_h.
\]

这与 prefill 中的 $S\times S$ attention 不同。Decode attention 不会物化完整 $S\times S$ matrix，因为只有一行 score：

\[
1\times S.
\]

因此 decode 的瓶颈更多是 KV Cache 读取和 batch 组织，而不是 $S^2$ 中间矩阵物化。FlashAttention 的 online softmax 和 fused kernel 思想仍然有帮助，但 decode kernel 通常需要针对 single-query 或 small-query 情况专门优化。

对于 batch decode，多个请求可以组成：

\[
B_{\mathrm{decode}}\times 1
\]

的 query batch，每个请求上下文长度不同。若配合 PagedAttention，K/V 还可能是分页存储。此时 attention kernel 需要同时处理：

\begin{itemize}
\item variable context length；
\item paged KV block table；
\item GQA/MQA；
\item small Q length；
\item memory-bound KV reads；
\item sampling 后的 dynamic scheduling。
\end{itemize}

因此，serving 系统中常常同时需要 FlashAttention 风格的 prefill kernel 和 PagedAttention 风格的 decode kernel。二者解决的问题不同：

\begin{itemize}
\item FlashAttention：减少 attention 计算中的 HBM IO；
\item PagedAttention：管理动态 KV Cache 并支持 continuous batching。
\end{itemize}

\begin{table}[H]
\centering
\small
\caption{Prefill 与 Decode 对 Attention Kernel 的不同要求}
\label{tab:prefill-decode-attention-kernel}
\begin{tabular}{p{0.2\textwidth}p{0.34\textwidth}p{0.34\textwidth}}
\toprule
\textbf{维度} & \textbf{Prefill} & \textbf{Decode} \
\midrule
Q 长度
& prompt 长度 $S$
& 通常为 1 \
\midrule
K/V 长度
& prompt 长度 $S$
& 历史上下文长度 $S_{\mathrm{cache}}$ \
\midrule
主要问题
& $S\times S$ attention IO 和计算
& KV Cache 读取、动态 batch、memory-bound \
\midrule
FlashAttention 价值
& 非常高，避免物化 attention matrix
& 思想有用，但通常需要 decode-specialized kernel \
\midrule
PagedAttention 价值
& 管理写入和长 prompt blocks
& 非常高，支持动态 KV 读取与 continuous batching \
\bottomrule
\end{tabular}
\end{table}

\subsection{与 PagedAttention 的互补关系}
\label{subsec:flashattention-pagedattention-complementary}

PagedAttention 和 FlashAttention 很容易被混淆，因为它们都与 Attention 和显存有关。但二者关注点不同。

\[
\text{FlashAttention}
\Rightarrow
\text{How to compute attention with less HBM IO?}
\]

\[
\text{PagedAttention}
\Rightarrow
\text{How to store and manage KV Cache for dynamic requests?}
\]

在一个现代 LLM serving engine 中，它们可能同时存在：

\begin{itemize}
\item prefill 使用 FlashAttention 或其变体；
\item decode 使用 PagedAttention kernel 读取分页 KV Cache；
\item 长 prompt chunked prefill 中也可能结合 paged KV 写入；
\item prefix cache 使用 block table 共享 KV；
\item attention backend 根据阶段、shape、GPU 架构选择不同 kernel。
\end{itemize}

可以把关系理解为：

\[
\text{FlashAttention 是计算优化，PagedAttention 是内存管理与调度优化。}
\]

\begin{figure}[H]
\centering
\begin{tikzpicture}[
font=\small,
box/.style={draw, rounded corners=4pt, minimum width=2.35cm, minimum height=0.8cm, align=center},
arrow/.style={-{Latex[length=2.1mm]}, thick}
]

\node[font=\bfseries] at (0,3.25) {FlashAttention 与 PagedAttention 在 Serving 中的互补};

\node[box, fill=green!10] (prefill) at (-4.2,1.35) {Prefill\\long prompt};
\node[box, fill=orange!14] (flash) at (-1.35,1.35) {FlashAttention\\IO-aware compute};
\node[box, fill=blue!8] (kvwrite) at (1.45,1.35) {Write KV\\blocks};
\node[box, fill=purple!10] (decode) at (4.2,1.35) {Decode\\token loop};

\node[box, fill=yellow!16] (paged) at (1.45,0.0) {PagedAttention\\paged KV read};
\node[box, fill=red!8] (sched) at (4.2,0.0) {Continuous\\batching};

\draw[arrow] (prefill) -- (flash);
\draw[arrow] (flash) -- (kvwrite);
\draw[arrow] (kvwrite) -- (decode);
\draw[arrow] (paged) -- (decode);
\draw[arrow] (sched) -- (decode);
\draw[arrow] (kvwrite) -- (paged);

\node[align=left, text width=10.5cm] at (0,-1.55) {
  Prefill 阶段常用 FlashAttention 减少 $S\times S$ attention 的 IO；
  Decode 阶段常用 PagedAttention 管理和读取动态 KV Cache。
};

\end{tikzpicture}
\caption{FlashAttention 与 PagedAttention 在推理系统中的互补关系}
\label{fig:flashattention-pagedattention-complementary}
\end{figure}

\section{常见误区与调试方法}
\label{sec:flashattention-misconceptions-debugging}

\subsection{误区一：FlashAttention 是近似 Attention}
\label{subsec:misconception-flashattention-approximate}

FlashAttention 不是近似 attention。它计算的是与标准 attention 数学上等价的精确结果，只是改变计算顺序和内存访问方式。由于浮点数加法和指数运算顺序变化，结果可能与标准实现存在极小数值差异，但这不是近似算法意义上的误差。

这点很重要。FlashAttention 不像某些 sparse attention、linear attention 或 low-rank attention 那样改变模型的 attention pattern。它的目标是在不改变模型语义的前提下提高速度和降低显存。

\subsection{误区二：FlashAttention 降低了 $O(S^2)$ 计算复杂度}
\label{subsec:misconception-flashattention-compute-complexity}

FlashAttention 不改变 dense attention 的 $O(S^2d_h)$ 计算复杂度。对于 full attention，所有 query-key 对仍然需要计算。它降低的是 HBM IO 和中间显存，不是数学计算阶数。

因此，当序列极长时，FlashAttention 仍然要面对 $S^2$ 计算量。若需要从复杂度上降低长上下文成本，需要考虑 sparse attention、sliding window attention、state space model、linear attention、attention sink、KV compression 等其他方向。

\subsection{误区三：所有场景都一定更快}
\label{subsec:misconception-flashattention-always-faster}

FlashAttention 在许多长序列和训练/prefill 场景中非常有效，但并不意味着所有 shape 都一定更快。影响因素包括：

\begin{itemize}
\item sequence length；
\item head dimension；
\item batch size；
\item causal 或 non-causal；
\item GPU 架构；
\item dtype；
\item dropout 是否开启；
\item mask 类型；
\item 是否 variable length；
\item kernel 是否支持当前 shape；
\item 是否 decode single-query。
\end{itemize}

对于很短序列，标准实现或其他 fused attention 可能已经足够快；对于某些特殊 mask 或 layout，FlashAttention kernel 可能无法直接使用；对于 decode 阶段，专门的 paged decode kernel 可能比通用 FlashAttention 更合适。

\subsection{调试与验证}
\label{subsec:flashattention-debug-validation}

接入 FlashAttention 或替换 attention backend 时，需要验证：

\begin{itemize}
\item 输出数值与标准实现是否在 tolerance 内一致；
\item backward 梯度是否正确；
\item causal mask 是否正确；
\item padding mask 或 variable length 是否正确；
\item dropout 随机性是否符合训练要求；
\item mixed precision 是否稳定；
\item 长序列下是否 OOM；
\item profiler 中 HBM traffic 是否下降；
\item kernel 是否真正被调用，而不是 fallback；
\item 不同 GPU 架构是否走了合适 kernel。
\end{itemize}

一个最小数值验证可以写成：

\begin{lstlisting}[language=Python,caption={FlashAttention 与标准 Attention 的数值验证示意},label={lst:flashattention-correctness-check}]
import torch
import torch.nn.functional as F

def standard_attention(q, k, v, causal=True):
# q, k, v: [batch, heads, seq, dim]
scale = q.shape[-1] ** -0.5
scores = torch.matmul(q, k.transpose(-1, -2)) * scale

if causal:
    seq = q.shape[-2]
    mask = torch.triu(
        torch.ones(seq, seq, device=q.device, dtype=torch.bool),
        diagonal=1,
    )
    scores = scores.masked_fill(mask, float("-inf"))

prob = torch.softmax(scores, dim=-1)
return torch.matmul(prob, v)

def check_attention_impl(q, k, v, flash_attention_fn):
ref = standard_attention(q, k, v, causal=True)
out = flash_attention_fn(q, k, v, causal=True)

max_abs_err = (ref - out).abs().max().item()
max_rel_err = ((ref - out).abs() / ref.abs().clamp_min(1e-6)).max().item()

print("max abs err:", max_abs_err)
print("max rel err:", max_rel_err)

torch.testing.assert_close(out, ref, rtol=1e-2, atol=1e-2)

\end{lstlisting}

在 BF16/FP16 下，不应要求逐 bit 一致。应使用合理的 absolute/relative tolerance，并分别验证不同 sequence length、head dimension、causal mask、batch size 和 dtype。

\section{本章小结}
\label{sec:chapter14-summary}

本章系统介绍了 FlashAttention 的动机、online softmax 推导、核心算法、前反向计算和工程实现视角。

\begin{itemize}
\item \textbf{标准 Attention 的主要问题之一是 IO 瓶颈}：它常常显式物化 $S\times S$ score matrix 和 softmax probability matrix，导致大量 HBM 读写。
\item \textbf{FlashAttention 不改变 Attention 数学结果}，而是通过 IO-aware tiling 和 kernel fusion 减少 HBM traffic。
\item \textbf{FlashAttention 不降低 dense attention 的 FLOPs 阶数}，仍然是 $O(S^2d_h)$，但显著降低中间显存和 HBM 访问。
\item \textbf{HBM 与片上 SRAM/register 的差异是理解 FlashAttention 的关键}：把中间 attention block 留在片上计算，通常比写入 HBM 再读回更快。
\item \textbf{Online softmax 使分块 attention 成为可能}：通过维护每行的运行最大值 $m$、指数和 $l$ 和输出累积 $a$，可以分块得到与完整 softmax 等价的结果。
\item \textbf{FlashAttention forward 按 Q/K/V blocks 计算}，局部 score、mask、softmax 和 $PV$ 累积都在 SRAM/register 中完成，最终只写回输出和少量统计量。
\item \textbf{FlashAttention backward 使用重计算换 IO}：不保存完整 attention probability，而是保存 $m,l$，在 backward 中分块重算需要的 softmax block。
\item \textbf{FlashAttention-2 进一步优化 work partitioning}，提升 occupancy，减少非矩阵乘法 FLOPs 和 shared memory 开销。
\item \textbf{FlashAttention-3 面向新硬件特性}，利用异步数据搬运、warp specialization、TMA 和 FP8 等能力继续提升性能。
\item \textbf{工程实现需要综合考虑 block size、shared memory、register pressure、warp 分工、Tensor Core 利用率和 mask 支持}。
\item \textbf{在推理系统中，FlashAttention 与 PagedAttention 是互补的}：前者主要优化 attention 计算 IO，后者主要优化动态 KV Cache 管理和 continuous batching。
\item \textbf{FlashAttention 不是所有 shape 都必然更快}。短序列、特殊 mask、single-query decode 或不支持的 layout 可能需要其他 attention backend。
\end{itemize}

下一章将进入量化技术。FlashAttention 通过减少 IO 提升 attention kernel 效率，而量化则从数据表示本身入手，降低模型权重、激活和 KV Cache 的字节数。我们会系统讨论 PTQ、QAT、GPTQ、AWQ、weight-only quantization、per-channel/per-group scale，以及量化对显存、带宽、吞吐和精度的影响。
